\chapter*{Abstract}
\addcontentsline{toc}{chapter}{Abstract}
\noindent
We describe an end-to-end pipeline classifying transient alerts as Type
Ia, Ib/c, or II supernovae, or flagging those resembling none, across
three data tiers: simulated ELAsTiCC2 light curves (\TierATotalObjects{}
objects, Tier A) for pretraining, real labelled ZTF light curves (BTSbot,
Tier B) for fine-tuning, and a live Fink stream (ZTF and LSST, Tier C) for
serving. The classifier is a causal gated recurrent unit over
per-observation flux, conditioned on continuous filter wavelength rather
than band identity, so weights learned on simulated LSST photometry
transfer to real ZTF photometry. A separate branch scores anomalies via
light-curve reconstruction error, Mahalanobis distance in the
classifier's latent space, and an isolation forest. On held-out Tier A
data the classifier reaches macro-F1 \TierAFullMacroF{} (Ia-vs-rest
ROC-AUC \LCOnlyRocAucIa{}); fine-tuned on Tier B, the deployed model
reaches \LCOnlyMacroF{} (a \SimToRealGapLCOnly{} drop, the expected
sim-to-real direction), and fine-tuning recovers \TransferGain{} macro-F1
over the zero-shot baseline. The anomaly branch's three-component
combination reaches ROC-AUC \AnomalyTierACombinedThreeRecomputed{} on
held-out simulated exotic classes and
\AnomalyTierBCombinedThreeRecomputed{} on real confirmed non-supernova
objects -- smaller but useful under domain shift. The pipeline was
validated against four pre-registered gates (reader correctness,
causality/label-shuffling, sim-to-real drop, manual anomaly inspection)
and briefly run against live ZTF and LSST alert traffic; a serving
latency benchmark exists in project records but is not restated here
since no artifact survives. We report two unanticipated results -- a
net-negative effect of host-galaxy features on real-data fine-tuning, and
an anomaly classifier that grows more, not less, confident on unfamiliar
inputs -- as open findings.
